\chapter{Intelligence and Analytical Methodology}
\label{ch:analytical-methodology}

This chapter details the mathematical, algorithmic, and analytical methodologies implemented in the Wolverine intelligence pipeline. It formalizes the probabilistic entity resolution scoring model, defines the string comparator and temporal proximity functions, documents the in-memory graph traversal architecture, and details the prompt-sanitized Retrieval-Augmented Generation (RAG) assistant.

\section{Analytical Pipeline Overview}
\label{sec:algo-pipeline}
The Wolverine analytical pipeline transforms raw, heterogeneous observations into correlated intelligence through a four-stage process:
\begin{enumerate}
    \item \textbf{Blocking Key Generation:} Partitions incoming account records into candidate comparison blocks to avoid quadratic $|A| \times |B|$ comparisons.
    \item \textbf{Heuristic Multi-Feature Scoring:} Evaluates pairwise similarity across alias strings, display names, temporal registration windows, baseline activity constants, and prefix cues.
    \item \textbf{Graph Projection:} Projects resolved identity clusters and observational interactions into an in-memory relational graph.
    \item \textbf{Sanitized RAG Synthesis:} Employs a local language model with strict prompt-injection guardrails to synthesize case narratives from retrieved evidence.
\end{enumerate}

\section{Entity Resolution and Blocking Strategy}
\label{sec:algo-entity-resolution}
To maintain low latency on single-host hardware, the entity resolver (\texttt{wolverine/src/resolver/entity\_resolver.ts}) generates three deterministic blocking keys for every account record:
\begin{itemize}
    \item \textbf{Handle Prefix Key:} First 3 characters of the normalized alphanumeric handle (e.g., \texttt{"sha"} from \texttt{"shadow\_runner"}).
    \item \textbf{Temporal Window Key:} ISO-8601 registration week string (e.g., \texttt{"2025-W03"}).
    \item \textbf{Domain Key:} Synthetic email provider domain (e.g., \texttt{"atlas.example.test"}).
\end{itemize}
Comparisons are strictly \textbf{cross-site only}; pairwise evaluations between accounts residing on the same site are pruned immediately ($S(e_a, e_b) = 0$).

\section{Mathematical Formulations}
\label{sec:algo-math}

\subsection{Jaro-Winkler String Similarity}
For alias and display name comparisons, the system implements the Jaro-Winkler metric~\cite{jaro1989advances,winkler1990string}.

Let $s_1$ and $s_2$ be two input strings. Two characters $c_1 \in s_1$ and $c_2 \in s_2$ are considered matching if they are identical and their positional distance satisfies:
\begin{equation}
\label{eq:jaro-window}
|pos(c_1) - pos(c_2)| \le \left\lfloor \frac{\max(|s_1|, |s_2|)}{2} \right\rfloor - 1
\end{equation}

Let $m$ be the number of matching characters and $t$ be the number of half-transpositions (matching characters that appear in different sequential order divided by 2). The Jaro similarity $s_j$ is given by:
\begin{equation}
\label{eq:jaro-similarity}
s_j(s_1, s_2) = \begin{cases} 
0 & \text{if } m = 0 \\
\frac{1}{3} \left( \frac{m}{|s_1|} + \frac{m}{|s_2|} + \frac{m - t}{m} \right) & \text{if } m > 0 
\end{cases}
\end{equation}

Winkler's modification adjusts $s_j$ based on the length $l$ of the common prefix (capped at $l \le 4$) using a constant prefix scaling factor $p = 0.1$:
\begin{equation}
\label{eq:jaro-winkler}
s_{jw}(s_1, s_2) = s_j(s_1, s_2) + l \cdot p \cdot (1 - s_j(s_1, s_2))
\end{equation}

\subsection{Temporal Proximity Decay}
The temporal proximity score evaluates whether two accounts were registered within a plausible temporal window:
\begin{equation}
\label{eq:temporal-proximity}
s_{\text{temp}}(t_a, t_b) = \max\left(0, 1 - \min\left(1, \frac{|t_a - t_b|}{30 \times 86400}\right)\right)
\end{equation}
Where $t_a$ and $t_b$ are Unix timestamps (seconds) of account creation. If the interval exceeds 30 days, $s_{\text{temp}}$ evaluates to 0.

\subsection{Weighted Multi-Feature Scoring Model}
The composite resolution score $S(e_a, e_b) \in [0.0, 1.0]$ between cross-site accounts $e_a$ and $e_b$ is defined by the weighted linear combination:
\begin{equation}
\label{eq:composite-score}
S(e_a, e_b) = w_a \cdot s_{\text{alias}} + w_d \cdot s_{\text{disp}} + w_t \cdot s_{\text{temp}} + w_o \cdot s_{\text{act}} + w_c \cdot s_{\text{cue}}
\end{equation}

The exact feature weights implemented in \texttt{entity\_resolver.ts} are:
\begin{itemize}
    \item $w_a = 0.30$: Alias similarity $s_{\text{alias}} = s_{jw}(e_a.\text{handle}, e_b.\text{handle})$
    \item $w_d = 0.15$: Display name similarity $s_{\text{disp}} = s_{jw}(e_a.\text{displayName}, e_b.\text{displayName})$
    \item $w_t = 0.20$: Temporal registration proximity $s_{\text{temp}}$
    \item $w_o = 0.15$: Activity overlap baseline $s_{\text{act}}$
    \item $w_c = 0.20$: Cross-site handle cue $s_{\text{cue}}$
\end{itemize}

\subsection{Documented Heuristic Simplifications}
Forensic code audit confirms two critical heuristic simplifications in the production prototype:
\begin{enumerate}
    \item \textbf{Hardcoded Activity Overlap ($s_{\text{act}}$):} In the current prototype (\texttt{entity\_resolver.ts} line 151), behavioral overlap is not computed from transactional logs. Instead, it is hardcoded to a static constant:
    \begin{equation}
    \label{eq:act-hardcode}
    s_{\text{act}} = 0.5 \quad (\text{contributing a static } +0.075 \text{ to all scores})
    \end{equation}
    \item \textbf{Binary Prefix Cross-Site Cue ($s_{\text{cue}}$):} Cross-site correlation does not perform Natural Language Processing (NLP) or stylometric analysis. It evaluates a binary 4-character prefix match:
    \begin{equation}
    \label{eq:cue-hardcode}
    s_{\text{cue}} = \begin{cases} 
    1.0 & \text{if } e_a.\text{handle}[0:4] == e_b.\text{handle}[0:4] \\
    0.0 & \text{otherwise}
    \end{cases}
    \end{equation}
\end{enumerate}

\subsection{Decision Logic and Thresholds}
Candidate pairs are classified according to three operational thresholds:
\begin{equation}
\label{eq:decision-thresholds}
\text{Decision}(S) = \begin{cases} 
\text{LINKED} & \text{if } S \ge 0.92 \quad (\text{Automatic identity assertion}) \\
\text{REVIEW} & \text{if } 0.70 \le S < 0.92 \quad (\text{Queued for analyst review}) \\
\text{REJECTED} & \text{if } S < 0.70 \quad (\text{Discarded candidate})
\end{cases}
\end{equation}

\section{In-Memory Graph Construction and Traversal}
\label{sec:algo-graph}
Earlier planning documents suggested the use of PostgreSQL recursive Common Table Expressions (CTEs) for graph traversal. Forensic inspection verifies that the operational engine (\texttt{wolverine/src/graph/projector.ts}) is implemented as a dedicated \textbf{In-Memory Adjacency Store}.

\subsection{Graph Representation}
The graph maintains nodes and directed edges using in-process hash maps:
\begin{itemize}
    \item \texttt{nodes}: $\text{Map}\langle \text{string}, \text{GraphNode} \rangle$
    \item \texttt{edges}: $\text{Map}\langle \text{string}, \text{GraphEdge} \rangle$
\end{itemize}
Edge types include \texttt{AUTHORED}, \texttt{REPLIED\_TO}, \texttt{LISTED}, \texttt{INTERACTED\_WITH}, and \texttt{POSSIBLE\_SAME\_AS} (which carries the confidence score $S$ from entity resolution).

\subsection{Breadth-First Search (BFS) Traversal}
Multi-hop queries execute an in-memory BFS queue with two strict runtime bounds:
\begin{itemize}
    \item Maximum traversal depth: $d_{\max} = 4$.
    \item Maximum returned records: $N_{\max} = 500$.
\end{itemize}
The computational complexity is $O(|V| + |E|)$ with negligible memory lookup latency, bounded by host RAM.

\section{Retrieval-Augmented Generation (RAG) and Guardrails}
\label{sec:algo-rag}
The AI Assistant (\texttt{wolverine/src/ai/rag.ts}) integrates local Ollama (Llama-3.2) inference to provide analytical question-answering over canonical records.

\subsection{4-Rule Input Sanitization}
To protect against adversarial prompt injection embedded in synthetic user content~\cite{perez2022ignore}, all retrieved context passes through four sequential filters:
\begin{enumerate}
    \item \textbf{Control Character Elimination:} Strips Unicode control characters (\texttt{U+0000}--\texttt{U+001F}, \texttt{U+007F}).
    \item \textbf{Length Truncation:} Enforces a strict 500-character maximum per field.
    \item \textbf{Pattern Stripping:} Removes adversarial prompt override tokens (e.g., \texttt{"ignore previous"}, \texttt{"system:"}, \texttt{"<|"}, \texttt{"|>"}, \texttt{"[INST]"}, \texttt{"[/INST]"}).
    \item \textbf{Token Bounding:} Bounds retrieved context to a maximum of 20 relevant records.
\end{enumerate}

\subsection{Deterministic Heuristic Fallback}
If Ollama is offline or fails to respond within a 5-second timeout window, the RAG engine transitions to a \textbf{Deterministic Fallback Mode}. It synthesizes a structured analytical hypothesis by extracting the top-3 most recent canonical records matching query keywords, ensuring the analyst workstation remains functional without external model dependencies.

\section{Synthetic Population Generation}
\label{sec:algo-synthetic}
The synthetic generator (\texttt{generator/src/generator/engine.py}) generates reproducible multi-platform populations using a seeded pseudorandom pipeline across 13 deterministic phases (generating canonical persons, assigning site accounts with a 35\% handle mutation rate, generating listings, orders, messages, forum replies, and injecting controlled evaluation scenarios).
